% ================================================================
\section{Theoretical Connections to Optimal Control Theory}
\label{sec:appendix_theory}

This appendix situates the zero-iteration controller within
a broader theoretical landscape, drawing connections to three
companion analyses: the tunnel--threshold--turbulence (TTT)
framework for iterative optimal control~\cite{zhang2026ttt},
the renormalization-group (RG) interpretation of the DARE
iteration~\cite{zhang2026ym}, and the adjoint--barrier--controllability
(ABC) analysis of constrained optimal control~\cite{zhang2026abc}.
These connections explain \emph{why} the barrier-augmented
zero-iteration controller cannot enforce absolute state constraints
and suggest concrete paths beyond this limitation.

\subsection{The Zero-Iteration Controller as a Seed}
\label{sec:seed}

The zero-iteration control law
$u = u^*(\balpha) - K(\balpha)\,\delta x$
is a \emph{seed} in the sense of~\cite{zhang2026ttt}:
a fixed-point controller at one operating point
(the reference trajectory), extended to nearby points
by first-order sensitivity.

TTT's main theorem establishes four equivalent conditions:
\begin{equation}\label{eq:ttt_equiv}
  \text{bounded } P_j
  \;\Longleftrightarrow\;
  \text{spectral gap} > 0
  \;\Longleftrightarrow\;
  \text{Lyapunov descent}
  \;\Longleftrightarrow\;
  \text{tunnel exists.}
\end{equation}
Our DARE solution $P_\infty$ is the bounded cost-to-go
guaranteed by the first condition.
The ``tunnel'' is the sublevel set
$\mathcal{T} = \{x : \frac{1}{2} x^\top P_\infty x \leq V_{\max}\}$
inside which the seed controller maintains
exponential contraction---the region of attraction
whose radius scales as
$r \sim \|P_\infty\|^{-1/2}$.

From the Yang--Mills perspective~\cite{zhang2026ym},
the DARE iteration $P_{j+1} = f(P_j)$ is a discrete
\emph{renormalization group flow}: each step integrates
out one ``scale'' of the control problem
(one stage of the finite-horizon cost).
Convergence of $P_j \to P_\infty$ is the analogue
of the mass gap in Yang--Mills theory---the spectral gap
between the largest and second-largest eigenvalue of the
closed-loop transition matrix.
The contraction rate
$\bar\rho(A_{\mathrm{cl}}) < 1$ (Proposition~\ref{prop:contraction})
is precisely this gap.

\subsection{Why Zero-Iteration Cannot Enforce Absolute Constraints}
\label{sec:why_barrier_fails}

The circle-in-square experiment
(Table~\ref{tab:circle}) reveals that the barrier-augmented
controller violates the square bound by the same 88\,mm as the
unconstrained analytic controller.
We now explain this result through three complementary lenses.

\textbf{Costate rotation (ABC).}
The gain $K(\balpha)$ is the costate $\nabla V$ evaluated at
the reference (equilibrium) point~\cite{zhang2026abc}.
The full costate \emph{field} $\nabla V(x)$ varies spatially:
near a constraint boundary, $\nabla V$ rotates to point
away from the boundary, encoding the directional information
needed to steer the trajectory into the feasible region.
A constant-gain approximation $K \approx \nabla V(x^*)$ cannot
capture this rotation, regardless of the barrier parameter $\mu$.
The barrier Hessian~\eqref{eq:barrier_hess_x} modifies the
\emph{magnitude} of the gain (increasing the effective state cost
near constraints) but not its \emph{direction}---and it is the
directional information that determines constraint satisfaction
when the reference lies outside the feasible region.

\textbf{UV--IR decoupling (Yang--Mills).}
In the language of~\cite{zhang2026ym}, the barrier modifies
the ``UV source term'' $Q_\mu$ (the error-state cost in the
DARE) but cannot create ``IR contraction''
(the iterative projection that enforces absolute constraints
over the prediction horizon).
The zero-iteration controller operates entirely in the UV
regime---local, single-point, based on the equilibrium
Hessian.
Hard constraint satisfaction is an IR property: it requires
the full cost-to-go landscape $V(x)$, not just its Hessian
$\nabla^2 V(x^*)$ at one point.
Both regimes must cooperate---UV provides the seed, IR provides
the tunnel.

\textbf{Seed vs.\ tunnel (TTT).}
The seed alone does not constitute a tunnel.
The tunnel is constructed by iterative refinement
(ADMM iterations, in our setting) that progressively
tightens the constraint satisfaction.
This is precisely what the ADMM variants achieve via per-step
horizon bounds (Section~\ref{sec:circle_square}): each
iteration projects the predicted trajectory toward the
feasible set, building the tunnel one iteration at a time.
The zero-iteration controller, by definition, performs
no such construction.

\subsection{Beyond Zero-Iteration: Path Integral Extensions}
\label{sec:extensions}

Three paths to constraint-aware real-time control, ordered by
computational budget:

\textbf{1.\ MPPI-seed hybrid}
($O(N_s \cdot n) \approx 1500$ ops).
Use $K(\balpha)$ as the base policy and sample $N_s \approx 10$
perturbations around the seed trajectory.
Weight each sample by the Boltzmann factor
$e^{-\beta \cdot J}$ where $J$ includes a constraint
penalty~\cite{williams2017}.
The weighted average creates a soft ``tunnel'' via ensemble
averaging: samples that violate constraints receive low weight.
This is deterministic-time (no convergence loop) and respects
the TTT structure: the seed provides the center of the sampling
distribution, while the Boltzmann weighting constructs a
statistical approximation to the tunnel boundary.

\textbf{2.\ Costate lookup table}
($O(n \cdot m)$ + table lookup $\approx 500$ ops).
Precompute the costate field
$\nabla V(x)$ at a grid of positions near constraint
boundaries.
Online, look up the nearest costate vector and correct
$K(\balpha)$ with the directional information.
This adds ABC's ``missing ingredient''---the costate
\emph{field} rather than a single-point costate---to the
zero-iteration controller.
The lookup table size scales with the number of constraint
boundaries, not the state dimension, keeping the approach
compact for box-constrained systems.

\textbf{3.\ Constraint-shaped terminal cost}
($O(n \cdot m) \approx 260$ ops, same budget as the
Analytic Template).
Instead of adding the barrier Hessian at the equilibrium,
design $P_\infty$ via a DARE whose terminal cost encodes
the constraint geometry along the predicted trajectory.
This requires a one-step prediction
$x_{k+1} = A_d x_k + B_d u_k$ before the gain lookup
but otherwise preserves the ${\sim}260$-operation datapath.
The key difference from the current barrier approach:
the terminal cost is evaluated along the trajectory,
not at the equilibrium---providing the ``IR information''
that the equilibrium-only barrier lacks.

These three extensions correspond to TTT's nested game
structure~\cite{zhang2026ttt}: the MPPI-seed hybrid
(outer game, sampling/exploration), costate lookup
(middle game, adjoint/directional information), and
constraint-shaped terminal cost (inner game, cost landscape
shaping).
The tunnel emerges from the cooperation of all three
mechanisms, not from any single one---mirroring the
experimental finding that no single zero-iteration mechanism
(barrier, clipping, or gain interpolation) suffices for
hard constraint satisfaction.
